\section{Discussion}
\label{sec:discussion}

The findings support a conditional account of efficient ontology reasoning.
Restricting a language can create opportunities for simpler computation, but
the benefit also depends on the arrangement of facts, the available
justifications, and the questions being asked. In the present experiments,
sharing surviving justifications was valuable when many individuals could
reuse the same reasoning. It became a cost when the attempted justification
did not exist. Similarly, spending more effort choosing an evaluation order
helped one demanding query, but imposed overhead on adverse cases. These
results identify conditions for improvement rather than a universally
preferable strategy.

The external comparisons make this distinction especially important. One
rule insertion was approximately six times faster than ZodiacEdge, whereas
one rule deletion was approximately forty-one times slower. Neither result
alone characterizes the relative quality of the two approaches. Together,
they show why an average across unlike operations can conceal a substantial
weakness. They also motivate a more discriminating question: which observable
properties of a change predict whether shared reasoning, selective
reconsideration, or broader recomputation will be most effective? The current
evidence does not yet establish a reliable general answer.

The Bach update discrepancy adds a distinct qualification. Successful
recomputation of a state does not establish that every path of changes leading
to it will preserve the same conclusions. The surviving family connection
supplies a small, independently understandable test of this requirement.
An observed failure in one recorded version warrants a reproducible
counterexample and a scoped conclusion, rather than a general ranking of
reasoning methods.

\subsection{The continuing value of restricted languages}

A fragment's scientific value does not depend on knowing how frequently it
is used in practice. A proven reduction in computational difficulty is
valuable even without an adoption survey. Prior work on description logics
whose statements take a restricted implication form, known as Horn description
logics, illustrates this point: restricting the form of logical statements
changes the established complexity of specified reasoning tasks
\cite{hustadt2005}. That result concerns precisely defined languages and
tasks. It does not, by itself, supply a complexity guarantee for the DLP L3
fragment or the reasoning procedure studied here.

Measured performance provides a different kind of evidence. On the earlier
constructed existential family, the specialized reasoner checked consistency and identified the
individuals satisfying the same questions faster than HermiT, which supports a broader OWL
language~\cite{glimm2014hermit}. This demonstrates a useful computational
specialization under the experimental conditions. It does not prove that all
DLP L3 tasks are easier, that the advantage grows indefinitely with input size,
or that restricting the language is the sole cause of the observed difference.
Much of the advantage was already present before the new optimizations.

That constructed family also belongs to OWL 2 EL, another established
restricted language. This qualification concerns the constructed experiment;
the full Bach ontology has a separate semantic scope. The result therefore leaves
open how the approach compares with reasoners specialized for EL. Moreover,
although the experimental statements require related individuals to exist,
a reasoner can derive the relevant class memberships without explicitly
constructing those individuals. The experiment establishes an advantage for the
specified questions; it does not establish an advantage of generating unnamed
individuals.

The relationship between the examined fragments and OWL 2 RL is one of
overlap, rather than a simple hierarchy. DLP L3 permits general statements
requiring the existence of a related individual, which lie outside RL's
permitted syntax. Conversely, RL includes constructs absent from the
reconstruction~\cite{owlprofiles2012}. It would therefore be misleading
to describe DLP L3 as a superset of RL or all the historical fragments as RL
subsets. Language membership, the meaning assigned to statements, and the
questions a procedure can answer completely must be assessed separately.

\subsection{Limits of the evidence and further questions}

The benchmark collection combines controlled examples, a university ontology,
the thesis's original Bach examples, changes to a published rule program,
and selected database query components.
This diversity tests several mechanisms, but does not constitute a
representative sample of ontology applications. The database components omit
parts of the full published workloads. Repeated measurements on one machine
provide local evidence, with limited information about variation across
hardware, systems, or larger datasets. Historical published timings provide
context, not directly comparable speed measurements.
The Bach examples retain their original small scale; repeating their
measurements does not turn them into a scalability study.

Agreement with independently justified answers strengthens confidence in the
reported experiments. It cannot establish correctness for every ontology,
nor certify the thesis's historical software. Likewise, corrections to
particular printed formulas and procedures do not invalidate the underlying
connection between ontology semantics and rule evaluation. They indicate
where the specification needs repair and where experimental reconstruction
can contribute to scientific scrutiny.

Further studies can test three propositions. First, the benefit of shared
justifications should increase with the proportion of individuals that share
both starting facts and surviving proofs. Second, selective rule-maintenance
methods should be most advantageous when a change affects a small part of
the reasoning dependencies. Third, choosing an evaluation order should help
only when avoided combinations compensate for the cost of making that choice.
Experiments that vary these properties independently would provide stronger
causal evidence than adding more undifferentiated benchmarks.

Broader comparisons should also include specialized EL reasoners and diverse
DLP L3 tasks, with explicit termination and completeness conditions. Some recursive
existential descriptions can make the present procedure generate an unending
sequence of unnamed individuals. Failure of this procedure to finish does not
establish that no other procedure could decide the same question. A separate
study could examine how often real ontologies and their application questions
fit each fragment. Such a survey would inform practical adoption, but its
difficulty does not postpone the theoretical or computational case for
studying restricted languages.
